\resizebox{0.88\textwidth}{!}{%
\begin{tikzpicture}[
    % Định dạng các khối theo sắc xám học thuật nhất quán với Hình 3.1
    rootbox/.style={draw=black, very thick, double, fill=gray!18, rounded corners=6pt, text width=7.2cm, align=left, inner sep=5pt, font=\small, text=black},
    cabox/.style={draw=black, thick, fill=gray!10, rounded corners=6pt, text width=7.2cm, align=left, inner sep=5pt, font=\small, text=black},
    leafbox/.style={draw=black, thick, fill=white, rounded corners=6pt, text width=7.2cm, align=left, inner sep=5pt, font=\small, text=black},
    storebox/.style={draw=black, thick, dashed, fill=gray!6, rounded corners=6pt, text width=6.6cm, align=center, inner sep=5pt, font=\small, text=black},
    % Kiểu mũi tên nét liền (tiến trình ký) và nét đứt (tiến trình thẩm định)
    issuancearrow/.style={-{Stealth[scale=1.05]}, thick, black},
    valarrow/.style={-{Stealth[scale=1.1]}, very thick, dash pattern=on 5pt off 2.5pt, black},
    trustarrow/.style={-{Stealth[scale=1.1]}, very thick, dash pattern=on 6pt off 3pt, black}
]

    % --- KHỐI GỐC TIN CẬY (TRUST ANCHOR TRÊN CÙNG) ---
    \node[storebox] (store) at (0, 15.2) {
        \textbf{Kho chứng thư tin cậy (Root Trust Store)}\\[1pt]
        {\scriptsize Quản lý bởi Hệ điều hành / Trình duyệt (Relying Party)\\[1pt]Chứa bản sao khóa công khai $PK_{\mathrm{Root}}$ được tin cậy mặc định}
    };

    % --- KHỐI 1: Root CA ---
    \node[rootbox] (root) at (0, 10.4) {
        \textbf{Mắt xích 1: Root CA (Trust Anchor)}\\[1pt]
        {\scriptsize \texttt{Subject: CN = Lab Root CA}}\\[1pt]
        {\scriptsize \texttt{Issuer:  CN = Lab Root CA} (Tự ký - Self-signed)}\\[1pt]
        {\scriptsize Khóa: Cặp khóa gốc $(SK_{\mathrm{Root}}, PK_{\mathrm{Root}})$}\\[1pt]
        {\scriptsize \texttt{Basic Constraints: critical, CA:TRUE}}\\[1pt]
        {\scriptsize \texttt{SKI:} Mã băm nhận diện khóa $PK_{\mathrm{Root}}$}
    };

    % --- KHỐI 2: Intermediate CA ---
    \node[cabox] (subca) at (0, 5.2) {
        \textbf{Mắt xích 2: Intermediate CA (Sub-CA)}\\[1pt]
        {\scriptsize \texttt{Subject: CN = Lab Intermediate CA}}\\[1pt]
        {\scriptsize \texttt{Issuer:  CN = Lab Root CA}}\\[1pt]
        {\scriptsize Khóa: Cặp khóa trung gian $(SK_{\mathrm{Sub}}, PK_{\mathrm{Sub}})$}\\[1pt]
        {\scriptsize \texttt{Basic Constraints: critical, CA:TRUE, pathlen:0}}\\[1pt]
        {\scriptsize \texttt{AKI:} Trỏ tới $\mathrm{SKI}_{\mathrm{Root}}$ \textbar{} \texttt{SKI:} Băm khóa $PK_{\mathrm{Sub}}$}
    };

    % --- KHỐI 3: Server Certificate ---
    \node[leafbox] (server) at (0, 0.0) {
        \textbf{Mắt xích 3: Server Certificate (Thực thể cuối)}\\[1pt]
        {\scriptsize \texttt{Subject: CN = server-a.test}}\\[1pt]
        {\scriptsize \texttt{Issuer:  CN = Lab Intermediate CA}}\\[1pt]
        {\scriptsize Khóa: Khóa công khai $PK_{\mathrm{Server}}$ ($SK_{\mathrm{Server}}$ lưu an toàn trên Server)}\\[1pt]
        {\scriptsize \texttt{Basic Constraints: CA:FALSE} \textbar{} \texttt{EKU: serverAuth}}\\[1pt]
        {\scriptsize \texttt{SAN: DNS:server-a.test} \textbar{} \texttt{AKI:} Trỏ tới $\mathrm{SKI}_{\mathrm{Sub}}$}
    };

    % Neo Root CA vào Trust Store
    \draw[trustarrow] (root.north) -- (store.south)
        node[midway, right=8pt, align=left, font=\footnotesize\bfseries] {Neo niềm tin\\[1pt](Trust Anchor)};

    % ==============================================================
    % TIÊU ĐỀ HAI TIẾN TRÌNH Ở ĐỈNH (KHÔNG XUNG ĐỘT)
    % ==============================================================
    \node[align=center, font=\footnotesize\bfseries, text=black] at (-6.3, 15.2) {
        [TIẾN TRÌNH KÝ PHÁT]\\[2pt]
        {\scriptsize (Dùng khóa bí mật $SK$)}
    };

    \node[align=center, font=\footnotesize\bfseries, text=black] at (6.5, 15.2) {
        [TIẾN TRÌNH THẨM ĐỊNH]\\[2pt]
        {\scriptsize (Dùng khóa công khai $PK$)}
    };

    % ==============================================================
    % TIẾN TRÌNH KÝ PHÁT (BÊN TRÁI - TỪ TRÊN XUỐNG)
    % ==============================================================
    % Root CA ký Sub-CA
    \draw[issuancearrow] (root.west) -- (-4.5, 10.4) -- (-4.5, 5.2) -- (subca.west);
    \node[anchor=east, align=right, font=\footnotesize] at (-4.7, 7.8) {
        \textbf{1. Ký phát Sub-CA:}\\[1pt]
        $\sigma_{\mathrm{Sub}} = \mathrm{Sign}_{SK_{\mathrm{Root}}}(\mathrm{TBS}_{\mathrm{Sub}})$\\[1pt]
        {\scriptsize $SK_{\mathrm{Root}}$ ký lên chứng thư Sub-CA}
    };

    % Sub-CA ký Server Certificate
    \draw[issuancearrow] (subca.west) -- (-4.5, 5.2) -- (-4.5, 0.0) -- (server.west);
    \node[anchor=east, align=right, font=\footnotesize] at (-4.7, 2.6) {
        \textbf{2. Ký phát Server:}\\[1pt]
        $\sigma_{\mathrm{Server}} = \mathrm{Sign}_{SK_{\mathrm{Sub}}}(\mathrm{TBS}_{\mathrm{Server}})$\\[1pt]
        {\scriptsize $SK_{\mathrm{Sub}}$ ký lên chứng thư Server}
    };

    % ==============================================================
    % TIẾN TRÌNH THẨM ĐỊNH RFC 5280 (BÊN PHẢI - TỪ DƯỚI LÊN)
    % ==============================================================
    % Thẩm định Server Certificate
    \draw[valarrow] (server.east) -- (4.7, 0.0) -- (4.7, 5.2) -- (subca.east);
    \node[anchor=west, align=left, font=\footnotesize] at (4.9, 2.6) {
        \textbf{3. Thẩm định Server Certificate:}\\[1pt]
        $\mathrm{Verify}_{PK_{\mathrm{Sub}}}(\sigma_{\mathrm{Server}}) = \mathrm{True}$\\[1pt]
        {\scriptsize $\bullet$ Đối chiếu: $\mathrm{Issuer} \equiv \mathrm{Subject}_{\mathrm{Sub}}$}\\[1pt]
        {\scriptsize $\bullet$ Khớp khóa: $\mathrm{AKI}_{\mathrm{Server}} \equiv \mathrm{SKI}_{\mathrm{Sub}}$}
    };

    % Thẩm định Intermediate CA
    \draw[valarrow] (subca.east) -- (4.7, 5.2) -- (4.7, 10.4) -- (root.east);
    \node[anchor=west, align=left, font=\footnotesize] at (4.9, 7.8) {
        \textbf{4. Thẩm định Intermediate CA:}\\[1pt]
        $\mathrm{Verify}_{PK_{\mathrm{Root}}}(\sigma_{\mathrm{Sub}}) = \mathrm{True}$\\[1pt]
        {\scriptsize $\bullet$ Đối chiếu: $\mathrm{Issuer} \equiv \mathrm{Subject}_{\mathrm{Root}}$}\\[1pt]
        {\scriptsize $\bullet$ Khớp khóa: $\mathrm{AKI}_{\mathrm{Sub}} \equiv \mathrm{SKI}_{\mathrm{Root}}$}
    };

\end{tikzpicture}%
}
